\documentclass[11pt]{article}
% Bien dich: pdflatex (hoac Overleaf). Hinh ve bang TikZ/pgfplots, khong can anh ngoai.
\usepackage[T5]{fontenc}
\usepackage[utf8]{inputenc}
\usepackage[vietnamese]{babel}
\usepackage[a4paper,margin=2.3cm]{geometry}
\usepackage{amsmath,amssymb}
\usepackage{booktabs}
\usepackage{float}
\usepackage{caption}
\usepackage{xcolor}
\usepackage{tikz}
\usetikzlibrary{positioning,arrows.meta,calc,fit,decorations.pathreplacing}
\usepackage{pgfplots}
\pgfplotsset{compat=1.17} 
\usepackage[hidelinks]{hyperref}

\definecolor{cOrig}{HTML}{7F7F7F}
\definecolor{cV2}{HTML}{D62728}
\definecolor{cH4}{HTML}{2CA02C}
\definecolor{cL2}{HTML}{1F77B4}
\definecolor{cBox}{HTML}{EAF2FB}
\definecolor{cNew}{HTML}{E8F5E9}
\definecolor{cWarn}{HTML}{FFF3E0}

\newcommand{\R}{\mathbb{R}}
\newcommand{\bmu}{\boldsymbol{\mu}}
\newcommand{\bSigma}{\boldsymbol{\Sigma}}
\newcommand{\vf}{\mathbf{f}}
\newcommand{\vx}{\mathbf{x}}

\title{Báo cáo dự án khoa học\\[4pt]
\large LBR v2: ``nhớ lại'' lớp cũ ở tầng $x'$ của khối cuối ViT\\
--- kết quả trên ImageNet-A (một bộ phân loại)}
\author{Thị Thanh Nga Hoàng}
\date{06/11/2026}

\begin{document}
\maketitle

\begin{abstract}
Các phương pháp không lưu ảnh (RSIAT, V2, HRC) ``nhớ'' mỗi lớp cũ bằng một đám mây Gauss đặt thẳng ở đặc trưng cuối
$\vf$, rồi huấn luyện lại bộ phân loại (CA) trên các điểm giả bốc từ đám mây đó. \textbf{LBR v2} lùi bộ nhớ này sớm hơn
một bước: lưu Gauss ở $x'$ --- token CLS ngay sau nhánh attention của khối Transformer cuối cùng --- và khi chạy CA thì
cho từng điểm giả đi qua \emph{phần còn lại của khối cuối hiện tại} (MLP, adapter, LayerNorm) để tạo đặc trưng giả. Như
vậy đặc trưng giả có đúng ``hình dáng'' phi tuyến do chính mạng tạo ra, và tự động theo kịp mọi thay đổi của adapter
khối cuối. Khác với LBR v1 (đã thất bại), bản v2 \emph{không} dùng điểm giả để huấn luyện adapter; nó chỉ thay đổi bước
CA. Trên ImageNet-A seed 1993, một lần chạy, LBR v2 đạt $\mathbf{75.32/66.03}$ (Avg/Last), cao hơn RSIAT gốc
($74.76/65.77$), V2 ($75.03/66.29$ --- nhưng Last thấp hơn) và LBR v1 tốt nhất ($75.25/65.77$); ngang HRC cùng số mẫu
$B=256$ ($75.33/65.83$) và thấp hơn HRC tốt nhất ($75.46/66.29$). Đây là bản LBR tốt nhất đến nay, nhưng chưa vượt HRC
và mới chỉ có một seed.
\end{abstract}

\tableofcontents
\newpage

% =====================================================================================================================
\section{Gốc rễ vấn đề, nói bằng lời thường}

\subsection{Bài toán}
Mô hình gồm \textbf{bộ trích đặc trưng} (ViT-B/16 huấn luyện sẵn trên ImageNet-21k, gắn thêm adapter nhỏ) biến ảnh
thành vector $\vf\in\R^{768}$, và \textbf{bộ phân loại} cosine gán ảnh cho lớp có trọng số ``cùng hướng'' nhất với
$\vf$. ImageNet-A có 200 lớp, chia thành 10 bước, mỗi bước 20 lớp mới, khoảng 30 ảnh huấn luyện mỗi lớp. Ở mỗi bước mô
hình chỉ thấy ảnh của 20 lớp mới và \textbf{không được giữ ảnh cũ}, nhưng vẫn phải phân loại đúng mọi lớp đã gặp.

\subsection{Cách ``nhớ lại'' hiện tại và chỗ yếu của nó}
Vì không còn ảnh, mỗi lớp cũ $c$ chỉ được nhớ bằng tâm $\bmu_c$ và hình dạng $\bSigma_c$ của đám đặc trưng. Sau mỗi
bước, CA sinh điểm giả từ $\mathcal{N}(\bmu_c,\bSigma_c)$ cho mọi lớp và huấn luyện lại bộ phân loại. Bộ phân loại chỉ
tốt bằng độ giống nhau giữa đám mây giả và đặc trưng thật. Có ba chỗ đám mây có thể sai:
\begin{enumerate}
\item \textbf{Sai vị trí} do adapter thay đổi làm đặc trưng trôi (drift). V2 xử lý bằng STS: ước lượng một ánh xạ
      tuyến tính giữa đặc trưng cũ và mới để dời $\bmu_c,\bSigma_c$.
\item \textbf{Sai độ rộng và hình dạng} vì $\bSigma_c$ ước lượng từ khoảng 30 ảnh trong 768 chiều. HRC xử lý bằng
      cách mượn hình dạng của lớp gần nhất và nới rộng đám mây.
\item \textbf{Sai ``kiểu'' phân phối}: đặc trưng cuối $\vf$ là đầu ra của khối MLP phi tuyến rồi qua LayerNorm, nên
      đám đặc trưng thật không phải một elip Gauss (ví dụ, LayerNorm ép mọi vector về gần cùng một độ dài). Đặt Gauss
      trực tiếp ở $\vf$ bỏ qua cấu trúc này. Hơn nữa, STS chỉ bù drift bằng một phép biến đổi \emph{tuyến tính}, trong
      khi thay đổi của khối cuối là phi tuyến.
\end{enumerate}
Ý (1) và (2) đã được V2 và HRC xử lý. \textbf{LBR v2 nhắm vào ý (3).} Lưu ý: ý (3) là giả thuyết động lực, chưa được đo
trực tiếp trong báo cáo này.

\subsection{Ý tưởng: lùi bộ nhớ về trước khối cuối một chút}
Khối Transformer cuối (AdaptFormer song song, như RSIAT) tính:
\begin{equation}
x'=x+\mathrm{Attn}(\mathrm{LN}_1(x)),\qquad
\vf=\mathrm{LN}\big(x'+\mathrm{MLP}(\mathrm{LN}_2(x'))+0.1\,\mathrm{Adapter}(x')\big)_{\text{CLS}}.
\end{equation}
Mọi phép tính sau Attention đều làm \emph{riêng trên từng token}. Vì vậy, chỉ cần token CLS $x'_{\text{CLS}}$ là tính
lại được \emph{chính xác} $\vf$ qua khối cuối hiện tại, không cần ảnh và không cần các token patch (kiểm tra số học:
sai số cỡ $6\cdot10^{-6}$). Do đó:
\begin{itemize}
\item \textbf{Lưu} Gauss $\mathcal{N}(\bmu'_c,\bSigma'_c)$ của $x'$ cho mỗi lớp, thay vì chỉ lưu Gauss ở $\vf$.
\item \textbf{Khi CA}: bốc $x'$ giả, cho đi qua MLP + adapter + LayerNorm của khối cuối \emph{hiện tại} để được $\vf$
      giả. Phần phi tuyến và phần drift của khối cuối do chính mạng tạo ra, không cần mô hình hoá.
\end{itemize}

\begin{figure}[H]
\centering
\begin{tikzpicture}[node distance=6mm and 7mm,
  blk/.style={draw,rounded corners,minimum height=8mm,align=center,font=\small,fill=cBox},
  mem/.style={draw,rounded corners,minimum height=8mm,align=center,font=\small,fill=cNew},
  ar/.style={-{Latex[length=2mm]},thick}]
\node[blk] (img) {Ảnh};
\node[blk,right=of img] (b1) {Khối 1..11\\(adapter)};
\node[blk,right=of b1] (att) {Attention\\khối 12};
\node[mem,right=of att,minimum width=12mm] (xp) {$x'_{\text{CLS}}$};
\node[blk,right=of xp] (mlp) {MLP + Adapter\\+ LayerNorm};
\node[blk,right=of mlp] (f) {$\vf$};
\draw[ar] (img)--(b1); \draw[ar] (b1)--(att); \draw[ar] (att)--(xp); \draw[ar] (xp)--(mlp); \draw[ar] (mlp)--(f);
\node[font=\footnotesize,align=center,below=9mm of xp,text=cL2] (lbrm) {LBR v2 lưu Gauss ở đây\\$\mathcal{N}(\bmu'_c,\bSigma'_c)$};
\node[font=\footnotesize,align=center,below=9mm of f,text=cV2] (v2m) {V2 / HRC lưu Gauss ở đây\\$\mathcal{N}(\bmu_c,\bSigma_c)$};
\draw[cL2,dashed,->] (lbrm)--(xp); \draw[cV2,dashed,->] (v2m)--(f);
\draw[decorate,decoration={brace,amplitude=4pt,mirror},thick] ($(mlp.south west)+(0,-0.15)$)--($(mlp.south east)+(0,-0.15)$);
\end{tikzpicture}
\caption{Vị trí của $x'$. Từ $x'_{\text{CLS}}$ đến $\vf$ chỉ là phép tính trên một token, nên có thể chạy lại trên
điểm giả.}
\end{figure}

% =====================================================================================================================
\section{Luồng LBR v2 trong một bước}\label{sec:flow}
Bước $t\ge1$ có $K$ lớp cũ và 20 lớp mới. Cấu hình chạy: \texttt{exps/lbr/imageneta\_lbr2\_obs\_xp\_s1993.json}
(\texttt{lbr\_lambda=0}, \texttt{lbr\_ca="xprime"}, \texttt{lbr\_ca\_kappa=true}).

\begin{enumerate}
\item \textbf{Huấn luyện adapter như V2.} CosFace trên lớp mới và RAE giữ đặc trưng mới gần cũ. Vì
      \texttt{lbr\_lambda}$=0$, LBR v2 \emph{không} thêm hàm mất nào vào bước này (chế độ ``quan sát''). Mọi tính toán
      phục vụ bộ nhớ $x'$ chạy với bộ sinh số ngẫu nhiên được khôi phục, nên không làm lệch quá trình huấn luyện.
\item \textbf{Dời bộ nhớ của lớp cũ bằng STS, chạy hai lần:} một lần ở $\vf$ (như V2) và một lần ở tầng $x'$. Cả
      hai dùng ảnh của bước hiện tại đi qua mô hình cũ và mô hình mới, rồi ước lượng ánh xạ (Procrustes) để dời
      $\bmu'_c,\bSigma'_c$ của $K$ lớp cũ. Ở tầng $x'$ không dùng mẫu đối kháng.
\item \textbf{Ghi nhớ lớp mới ở tầng $x'$.} Với mỗi lớp mới, cho ảnh huấn luyện đi qua mạng để lấy $x'$. Tính tâm
      $\bmu'_c$. Hình dạng $\bSigma'_c$ mượn theo kiểu HRC ($n_0=20$, $\tau=256$): trộn covariance của lớp với hình
      dạng của lớp gần nhất (cả cũ lẫn mới).
\item \textbf{CA ở tầng $x'$} (10 epoch, $B=256$ điểm mỗi lớp mỗi epoch):
  \begin{enumerate}
  \item bốc $x'\sim\mathcal{N}(\bmu'_c,\bSigma'_c)$ cho \emph{mọi} lớp, cũ lẫn mới;
  \item cho đi qua MLP + adapter + LayerNorm của khối cuối \emph{hiện tại} (không gradient) để được $\vf$ giả;
  \item dời tâm theo hệ số $\kappa_c=0.9+0.1(t_c+1)/(t+1)$ như V2 (cộng $(\kappa_c-1)\cdot$ tâm mẫu của lớp);
  \item chuẩn hoá và huấn luyện lại bộ phân loại cosine (thang 30, SGD + cosine) như CA thường.
  \end{enumerate}
\item \textbf{Đánh giá} trên toàn bộ tập kiểm tra của các lớp đã gặp. Bộ phân loại sau CA được dùng tiếp ở bước sau.
\end{enumerate}

\begin{figure}[H]
\centering
\begin{tikzpicture}[node distance=5mm,
  st/.style={draw,rounded corners,text width=3.0cm,minimum height=11mm,align=center,font=\small,fill=cBox},
  nw/.style={st,fill=cNew},
  ar/.style={-{Latex[length=2mm]},thick}]
\node[st] (s1) {1. Huấn luyện adapter\\(như V2, không thêm loss)};
\node[st,right=of s1] (s2) {2. STS: dời bộ nhớ\\lớp cũ ở $\vf$ và $x'$};
\node[nw,right=of s2] (s3) {3. Gauss $x'$ cho\\lớp mới (mượn hình)};
\node[nw,below=8mm of s3] (s4) {4a. Bốc $x'$ giả\\mọi lớp};
\node[nw,left=of s4] (s5) {4b. Qua khối cuối\\hiện tại $\Rightarrow\vf$ giả};
\node[st,left=of s5] (s6) {4c--d. $\kappa$, huấn luyện\\lại head cosine};
\draw[ar] (s1)--(s2); \draw[ar] (s2)--(s3); \draw[ar] (s3)--(s4); \draw[ar] (s4)--(s5); \draw[ar] (s5)--(s6);
\node[st,below=8mm of s6,fill=white] (s7) {5. Đánh giá;\\head dùng ở bước sau};
\draw[ar] (s6)--(s7);
\end{tikzpicture}
\caption{Luồng một bước. Ô màu xanh lá là phần mới của LBR v2. Phần còn lại giống V2.}
\end{figure}

\paragraph{Chi phí.} Mỗi epoch CA đẩy $200\times256=51\,200$ điểm qua MLP của khối cuối, chia theo cụm để tiết kiệm bộ
nhớ, và cần phân rã Cholesky cho 200 ma trận $768\times768$. Ở bước cuối, CA mất khoảng 7 giây. Bộ nhớ lưu thêm: một
mean và một covariance ở tầng $x'$ cho mỗi lớp, tức gấp đôi bộ nhớ thống kê của V2.

% =====================================================================================================================
\section{Điểm mới}

\begin{table}[H]
\centering\small
\caption{So sánh các phương pháp ở ba câu hỏi: lưu gì, đặc trưng giả tạo ra thế nào, adapter có học từ điểm giả không.}
\begin{tabular}{p{2.3cm}p{3.6cm}p{4.6cm}p{3.6cm}}
\toprule
& Bộ nhớ lớp cũ & Đặc trưng giả cho CA & Adapter học từ điểm giả?\\
\midrule
RSIAT gốc & Gauss ở $\vf$, dời tâm bằng SDC & Gauss ở $\vf$ & Không\\
V2 & Gauss ở $\vf$, dời bằng STS & Gauss ở $\vf$, $\kappa$ & Không\\
HRC & như V2, mượn hình dạng & Gauss ở $\vf$, mượn hình, $s=2$, $B=1024$ & Không\\
LBR v1 & Gauss ở $x'$ & (CA như V2) & \textbf{Có}: CE trên $\vf$ giả, gradient vào adapter khối cuối\\
\textbf{LBR v2} & Gauss ở $x'$ (+ ở $\vf$), dời bằng STS ở cả hai tầng & \textbf{$x'$ giả qua khối cuối hiện tại}: phi tuyến, khớp adapter hiện tại & Không\\
\bottomrule
\end{tabular}
\end{table}

\noindent Cụ thể, LBR v2 có ba điểm mới:
\begin{enumerate}
\item \textbf{Đặc trưng giả không còn là Gauss.} Đây là phương pháp đầu tiên trong chuỗi thí nghiệm mà đặc trưng giả
      được \emph{mạng thật} sinh ra từ một Gauss ở tầng sớm hơn, thay vì được giả định là Gauss ở tầng cuối.
\item \textbf{Bù drift của khối cuối chính xác, không cần ước lượng.} Thay đổi của adapter khối cuối được phản ánh
      trực tiếp trong $\vf$ giả. STS chỉ còn phải bù drift của khối 1--11 (ở tầng $x'$).
\item \textbf{Tách bộ nhớ $x'$ khỏi việc huấn luyện adapter.} Đây là bài học từ LBR v1: dùng điểm giả để huấn luyện
      adapter không có ích (adapter khối cuối gần như không đổi trong một bước) và còn làm hại lớp cũ. V2 giữ ý tưởng
      tái tạo ở $x'$ nhưng chỉ dùng nó cho CA, là nơi đám mây quyết định trực tiếp biên phân loại.
\end{enumerate}

% =====================================================================================================================
\section{Kết quả}
\paragraph{Thiết lập.} ImageNet-A, 10 bước $\times$ 20 lớp, ViT-B/16-IN21k + adapter, seed 1993, một bộ phân loại CA.
Avg là trung bình 10 bước; Last là độ chính xác sau bước cuối. Một ảnh kiểm tra tương đương $0.06$--$0.3$ điểm.
Log: \texttt{logs/adapter/imageneta/0/20/lbr2\_obs\_xp\_s1993\_*.log}.

\begin{table}[H]
\centering\small
\caption{ImageNet-A, seed 1993, độ chính xác top-1 (\%) sau CA ở từng bước.}
\setlength{\tabcolsep}{3.0pt}
\begin{tabular}{l*{10}{c}cc}
\toprule
Bước & 0 & 1 & 2 & 3 & 4 & 5 & 6 & 7 & 8 & 9 & Avg & Last\\
\midrule
RSIAT gốc & 86.29 & 82.50 & 79.62 & 78.65 & 75.51 & 73.23 & 70.56 & 68.59 & 66.88 & 65.77 & 74.76 & 65.77\\
V2 (chỉ CA) & 86.29 & 82.22 & 79.62 & 78.51 & 75.62 & 74.05 & 71.19 & 69.07 & 67.46 & \textbf{66.29} & 75.03 & \textbf{66.29}\\
LBR v1 tốt nhất & 86.29 & \textbf{84.44} & \textbf{81.09} & 78.37 & 75.51 & 73.95 & 71.45 & 68.51 & 67.10 & 65.77 & 75.25 & 65.77\\
HRC, $B=256$ & 86.29 & 83.89 & 80.88 & 78.37 & \textbf{75.86} & 73.64 & \textbf{71.81} & 69.23 & 67.46 & 65.83 & 75.33 & 65.83\\
HRC, $B=1024$ & 86.29 & 84.17 & 80.88 & 78.51 & \textbf{75.86} & 73.85 & 71.54 & 69.47 & \textbf{67.74} & \textbf{66.29} & \textbf{75.46} & \textbf{66.29}\\
\textbf{LBR v2} & 86.29 & 83.61 & 80.04 & \textbf{79.23} & 75.62 & \textbf{74.15} & 71.36 & \textbf{69.63} & 67.24 & 66.03 & 75.32 & 66.03\\
\midrule
LBR v2 $-$ V2 & 0.00 & $+1.39$ & $+0.42$ & $+0.72$ & 0.00 & $+0.10$ & $+0.17$ & $+0.56$ & $-0.22$ & $-0.26$ & $+0.29$ & $-0.26$\\
LBR v2 $-$ HRC$_{1024}$ & 0.00 & $-0.56$ & $-0.84$ & $+0.72$ & $-0.24$ & $+0.30$ & $-0.18$ & $+0.16$ & $-0.50$ & $-0.26$ & $-0.14$ & $-0.26$\\
LBR v2 $-$ gốc & 0.00 & $+1.11$ & $+0.42$ & $+0.58$ & $+0.11$ & $+0.92$ & $+0.80$ & $+1.04$ & $+0.36$ & $+0.26$ & $+0.56$ & $+0.26$\\
\bottomrule
\end{tabular}
\end{table}
LBR v1 tốt nhất là biến thể \texttt{lbr\_sa\_hrc\_bal\_lr3}. Bước 0 giống nhau ở mọi phương pháp vì chưa có CA.

\begin{figure}[H]
\centering
\begin{tikzpicture}
\begin{axis}[width=0.92\linewidth,height=6.2cm,xlabel={Bước},ylabel={Độ chính xác (\%)},xtick={0,...,9},
 ymin=64,ymax=88,grid=major,grid style={black!12},legend style={at={(0.02,0.03)},anchor=south west,font=\small}]
\addplot[cOrig,thick,mark=o] coordinates {(0,86.29)(1,82.50)(2,79.62)(3,78.65)(4,75.51)(5,73.23)(6,70.56)(7,68.59)(8,66.88)(9,65.77)};
\addlegendentry{RSIAT gốc}
\addplot[cV2,thick,mark=square] coordinates {(0,86.29)(1,82.22)(2,79.62)(3,78.51)(4,75.62)(5,74.05)(6,71.19)(7,69.07)(8,67.46)(9,66.29)};
\addlegendentry{V2 (chỉ CA)}
\addplot[cH4,thick,mark=*] coordinates {(0,86.29)(1,84.17)(2,80.88)(3,78.51)(4,75.86)(5,73.85)(6,71.54)(7,69.47)(8,67.74)(9,66.29)};
\addlegendentry{HRC, $B=1024$}
\addplot[cL2,very thick,mark=triangle*] coordinates {(0,86.29)(1,83.61)(2,80.04)(3,79.23)(4,75.62)(5,74.15)(6,71.36)(7,69.63)(8,67.24)(9,66.03)};
\addlegendentry{LBR v2}
\end{axis}
\end{tikzpicture}
\caption{Đường độ chính xác theo bước.}
\end{figure}

\begin{figure}[H]
\centering
\begin{tikzpicture}
\begin{axis}[ybar,bar width=6pt,width=0.92\linewidth,height=5cm,xlabel={Bước},ylabel={Chênh lệch (điểm)},
 xtick={1,...,9},ymin=-1,ymax=1.6,grid=major,grid style={black!12},
 legend style={at={(0.98,0.97)},anchor=north east,font=\small}]
\addplot[fill=cV2!60] coordinates {(1,1.39)(2,0.42)(3,0.72)(4,0.00)(5,0.10)(6,0.17)(7,0.56)(8,-0.22)(9,-0.26)};
\addlegendentry{LBR v2 $-$ V2}
\addplot[fill=cH4!60] coordinates {(1,-0.56)(2,-0.84)(3,0.72)(4,-0.24)(5,0.30)(6,-0.18)(7,0.16)(8,-0.50)(9,-0.26)};
\addlegendentry{LBR v2 $-$ HRC$_{1024}$}
\end{axis}
\end{tikzpicture}
\caption{LBR v2 hơn V2 ở bước 1--7 rồi thua nhẹ ở bước 8--9. So với HRC thì lúc hơn lúc kém.}
\end{figure}

\subsection{Lớp cũ và lớp mới ở bước cuối}
\begin{table}[H]
\centering\small
\caption{Bước 9: độ chính xác trên 180 lớp cũ và 20 lớp mới.}
\begin{tabular}{lccc}
\toprule
Phương pháp & Tổng (Last) & Lớp cũ & Lớp mới\\
\midrule
RSIAT gốc & 65.77 & 65.52 & 68.55\\
V2 (chỉ CA) & \textbf{66.29} & \textbf{66.16} & 67.74\\
LBR v1 tốt nhất & 65.77 & 65.30 & \textbf{70.97}\\
HRC, $B=256$ & 65.83 & 65.59 & 68.55\\
HRC, $B=1024$ & \textbf{66.29} & 66.02 & 69.35\\
\textbf{LBR v2} & 66.03 & 65.59 & \textbf{70.97}\\
\bottomrule
\end{tabular}
\end{table}

\subsection{CA đóng góp bao nhiêu ở mỗi bước}
Bảng dưới so sánh độ chính xác trước CA (head sau khi huấn luyện adapter) và sau CA. Lưu ý: vì bộ phân loại sau CA
được dùng tiếp ở bước sau, quỹ đạo của adapter ở ba lần chạy \emph{không} trùng nhau từ bước 2 trở đi.

\begin{table}[H]
\centering\small
\caption{Độ chính xác trước CA $\to$ mức tăng nhờ CA (điểm).}
\setlength{\tabcolsep}{3.4pt}
\begin{tabular}{l*{9}{c}cc}
\toprule
Bước & 1 & 2 & 3 & 4 & 5 & 6 & 7 & 8 & 9 & TB trước CA & TB tăng\\
\midrule
V2 & 82.22 & 79.62 & 75.07 & 74.20 & 72.41 & 69.94 & 67.39 & 66.88 & 64.78 & 72.50 & \\
\quad tăng nhờ CA & 0.00 & 0.00 & $+3.44$ & $+1.42$ & $+1.64$ & $+1.25$ & $+1.68$ & $+0.58$ & $+1.51$ & & $+1.28$\\
HRC$_{1024}$ & 82.22 & 80.25 & 76.07 & 75.15 & 71.38 & 69.58 & 67.39 & 66.81 & 64.91 & 72.42 & \\
\quad tăng nhờ CA & $+1.95$ & $+0.63$ & $+2.44$ & $+0.71$ & $+2.47$ & $+1.96$ & $+2.08$ & $+0.93$ & $+1.38$ & & $\mathbf{+1.62}$\\
LBR v2 & 82.22 & 80.25 & 76.07 & 75.51 & 72.10 & 69.76 & 67.71 & 67.10 & 64.78 & \textbf{72.83} & \\
\quad tăng nhờ CA & $+1.39$ & $-0.21$ & $+3.16$ & $+0.11$ & $+2.05$ & $+1.60$ & $+1.92$ & $+0.14$ & $+1.25$ & & $+1.27$\\
\bottomrule
\end{tabular}
\end{table}

% =====================================================================================================================
\section{Phân tích}
\begin{enumerate}
\item \textbf{Đây là bản LBR tốt nhất.} So với LBR v1 tốt nhất: Avg $+0.07$, Last $+0.26$. Bỏ việc huấn luyện adapter
      bằng điểm giả là đúng hướng. LBR v1 hơn ở bước 1--2 nhưng lớp cũ tụt dần; LBR v2 không bị tụt như vậy ở cuối
      chuỗi (lớp cũ ở bước 9: $65.59$ so với $65.30$).
\item \textbf{So với V2: Avg tăng, Last giảm nhẹ.} Phần tăng nằm ở bước 1--7 ($+1.39$ ở bước 1, $+0.72$ ở bước 3,
      $+0.56$ ở bước 7). Ở bước 8--9 thì thua $0.2$--$0.3$ điểm, tức khoảng 2--4 ảnh, nằm trong vùng nhiễu.
\item \textbf{So với HRC: ngang ở cùng điều kiện, thua bản tốt nhất.} Cùng $B=256$, LBR v2 ngang HRC
      ($75.32$ so với $75.33$ Avg) và hơn ở Last ($66.03$ so với $65.83$). HRC tốt nhất dùng thêm $s=2$ và $B=1024$,
      hai thứ LBR v2 \emph{chưa} được tune. Vì vậy so sánh hiện tại chưa công bằng hoàn toàn với LBR v2.
\item \textbf{Đổi lỗi ``mới bị đoán thành cũ'' lấy lỗi ``cũ bị đoán thành mới''.} Ở bước cuối, lớp mới đạt $70.97$
      (cao nhất, bằng LBR v1), so với $67.74$ của V2 và $69.35$ của HRC. Đổi lại, lớp cũ giảm từ $66.16$ (V2) xuống
      $65.59$. Báo cáo trước chỉ ra lỗi chính là lớp mới bị đoán thành lớp cũ, và LBR v2 giảm đúng loại lỗi đó, nhưng
      bù trừ lại khiến tổng gần như không đổi.
\item \textbf{Mức tăng do CA tự thân không lớn hơn V2.} Trung bình mỗi bước, CA của LBR v2 thêm $+1.27$ điểm, xấp xỉ V2
      ($+1.28$) và thấp hơn HRC ($+1.62$). Lợi thế Avg của LBR v2 chủ yếu đến từ việc độ chính xác \emph{trước CA}
      cao hơn ($72.83$ so với $72.50$/$72.42$): bộ phân loại sau CA ở bước trước ảnh hưởng tới việc huấn luyện adapter
      ở bước sau. Chưa thể nói hiệu ứng quỹ đạo này là ổn định hay chỉ là may rủi của một seed.
\item \textbf{CA ở tầng $x'$ dao động mạnh giữa các bước.} Ở bước 2, 4, 8, CA gần như không tăng ($-0.21$, $+0.11$,
      $+0.14$), trong khi HRC vẫn tăng $0.6$--$0.9$. Một giả thuyết: đám mây $x'$ chỉ có $s=1$ và $B=256$, nên sau khi
      qua LayerNorm vẫn còn hẹp hoặc phủ chưa đủ, đúng hiện tượng HRC đã gặp với đám mây hẹp của V2.
\end{enumerate}

% =====================================================================================================================
\section{Giới hạn}
\begin{itemize}
\item \textbf{Một seed, một lần chạy.} Theo đo trước đây, nhiễu giữa các seed khoảng $\pm1$ điểm Last, còn nhiễu của CA
      khoảng $\pm0.3$--$0.5$. Mọi chênh lệch dưới $0.5$ trong báo cáo này chưa đủ để kết luận.
\item \textbf{Chưa tune.} LBR v2 dùng nguyên $n_0=20$, $\tau=256$ của HRC (vốn tune cho tầng $\vf$), cùng với $s=1$,
      $B=256$. Siêu tham số của các phương pháp so sánh được tune trên tập kiểm tra seed 1993.
\item \textbf{Giả thuyết ``đặc trưng không Gauss'' chưa được đo trực tiếp}, ví dụ bằng cách so sánh phân phối $\vf$
      giả với $\vf$ thật của lớp cũ ở bước cuối.
\item \textbf{Quá trình chạy.} Trước lần chạy hoàn chỉnh, có các lần bị dừng do tràn bộ nhớ GPU và một lỗi làm mô hình
      cũ bị chuyển sang chế độ train khi tính STS ở tầng $x'$ (đã sửa trong \texttt{XPrimeNet}). Kết quả ở đây chỉ lấy từ lần chạy
      hoàn chỉnh; log của các lần chạy lỗi đã được xoá.
\end{itemize}

% =====================================================================================================================
\section{Kết luận và bước tiếp theo}
LBR v2 chuyển ý tưởng ``nhớ lại ở tầng $x'$'' từ việc huấn luyện adapter (v1, thất bại) sang bước CA. Ở đó, đặc trưng
giả được chính khối cuối hiện tại tạo ra. Kết quả $75.32/66.03$ là bản LBR tốt nhất, hơn RSIAT gốc $+0.56/+0.26$ và
hơn V2 $+0.29$ Avg, ngang HRC ở cùng $B=256$, nhưng chưa vượt HRC tốt nhất ($75.46/66.29$).

\paragraph{Đề xuất, theo thứ tự ưu tiên.}
\begin{enumerate}
\item \textbf{Áp dụng hai thứ đã có ích cho HRC} vào CA tầng $x'$: nới rộng $s=2$ và $B=1024$. Đây là phép thử rẻ
      nhất để biết LBR v2 có vượt được HRC hay không.
\item \textbf{Chạy thêm seed 1994, 1995} cho LBR v2, V2 và HRC để biết chênh lệch có thật hay không.
\item \textbf{Đo trực tiếp giả thuyết:} so sánh $\vf$ giả (qua khối cuối) với $\vf$ thật của lớp cũ, và với Gauss ở
      $\vf$ của V2/HRC, bằng khoảng cách phân phối hoặc oracle trên lớp cũ.
\item \textbf{Kết hợp với LDA} lúc đánh giá (\texttt{eval\_ens\_lda}). Hướng này từng tăng Last một cách nhất quán, và
      có thể bù phần lớp cũ bị giảm.
\end{enumerate}

\paragraph{Mã và cấu hình.} \texttt{models/lbr.py} (\texttt{xprime}, \texttt{head\_from\_xprime},
\texttt{train\_ca\_push}, \texttt{XPrimeNet}); \texttt{models/RSIAT\_adapter.py} (\texttt{\_ca\_xprime},
\texttt{\_lbr\_class\_stats}); \texttt{models/hrc.py} (\texttt{fixed\_shapes}); cấu hình
\texttt{exps/lbr/imageneta\_lbr2\_obs\_xp\_s1993.json}.

\end{document}
